\documentclass[10pt,a4paper]{article}
\usepackage[a4paper,margin=1.8cm]{geometry}
\usepackage[T1]{fontenc}
\usepackage{lmodern}
\usepackage{microtype}
\usepackage{booktabs}
\usepackage{tabularx}
\usepackage{array}
\usepackage[table]{xcolor}
\usepackage{enumitem}
\usepackage{titlesec}
\usepackage{url}
\usepackage[colorlinks=true,linkcolor=blue!50!black,urlcolor=blue!50!black]{hyperref}

\setlist{nosep,leftmargin=1.3em}
\titlespacing*{\section}{0pt}{1.1ex}{0.6ex}
\titlespacing*{\subsection}{0pt}{0.9ex}{0.4ex}
\setlength{\parindent}{0pt}
\setlength{\parskip}{0.45ex}
\renewcommand{\arraystretch}{1.08}
\newcommand{\iss}[1]{\href{https://github.com/JieGH/RVV_code_gen_via_MLIR_xDSL/issues/#1}{\##1}}
\newcommand{\up}[1]{\textcolor{red!70!black}{#1}}
\newcommand{\down}[1]{\textcolor{blue!60!black}{#1}}
\newcommand{\good}[1]{\textcolor{green!45!black}{\textbf{#1}}}

\title{\vspace{-1.2cm}\textbf{B3 and the microkernel tile:\\ the slow state follows the width of the B panel}\\[2pt]
\large Brief of issue \iss{48} (25 tiles, two packers, two boards), a follow-up of \iss{41}, \iss{44}, \iss{46} and \iss{47}}
\author{RVV\_code\_gen\_via\_MLIR\_xDSL, \texttt{pure-xdsl}\\ code \texttt{d7ce90e}, \texttt{ac6a43f}, \texttt{8d7562c}; logs \texttt{476d843}}
\date{1 October 2026}

\begin{document}
\maketitle
\vspace{-0.8cm}

\paragraph{Scope.} B3 (64$\times$384$\times$384) has two speeds: in most reps it falls into a slow state
(\iss{36}, \iss{41}, \iss{42}). Two earlier findings left the cause open. \iss{46} found that the reference design's B3
on a 16$\times$14 tile is steady while ours on 16$\times$15 is slow in 87--89\% of its reps; \iss{47} found that every
shape in the slow state had run on the 16$\times$15 tile and every steady large shape on 32$\times$6 or 32$\times$7, so
the tile and the number of rows could not be told apart. \iss{48} therefore asks one question: \emph{does the
microkernel tile (MR$\times$NR) decide whether the slow state appears, and does the packing side of B3 depend on it?}
This brief reports a sweep of 25 tiles (MR from 4 to 32, NR from 4 to 31), each run on four shapes (B3, B3 with
K=256, B2 and Reuse24), with the vector packer and, on 7 tiles, the scalar packer. Both boards ran at the same time,
one series each, pinned to core 4. \textbf{Every speed number is from the uninstrumented build}
(\texttt{INSTRUMENTED=0}), reference protocol, \texttt{-{}-tmin 10}, B3 at the reference's nc=384
(\texttt{-{}-no-nc-overrides}); speed is in GOPs, by mean time (the reference's formula) and by median rep. The
instrumented build was used only to split each rep into pack time and compute time. Launches per case:
\texttt{jbpi2} 4 (vector and scalar) and 3 (instrumented); \texttt{jbpi} 3 throughout. Every log cited is committed
(Appendix~A).

\section{Findings}
\begin{enumerate}
  \item \textbf{The slow state follows NR.} With the vector packer, B3 is \good{steady} (no slow rep in any launch, on
  either board) on every tile with NR of 4 or 6, and \up{two-speed} (53--94\% slow reps; 16$\times$16, over the register budget, 25--67\%) on every tile with NR of 7 or
  more. This holds from MR\,=\,4 to MR\,=\,32.
  \item \textbf{It does not follow MR or kc.} Steady and slow tiles share the same MR values, and the cache model's kc
  (128, 170, 256 or 384) is present among both.
  \item \textbf{The scalar packer is far steadier from NR\,=\,8 upwards.} On 16$\times$8 it runs at about 7.0--7.3 GOPs
  with 2--8\% slow reps; the vector packer runs at 3.6--3.7 GOPs with 89--94\% slow reps. At NR of 6 or less the
  vector packer is level with or ahead of the scalar one.
  \item \textbf{In a slow rep both halves slow down:} packing takes 1.6--2.2 times as long and compute 1.2--1.8 times.
  \item \textbf{Rows and depth matter as well (\iss{47}).} 32$\times$7 is steady on the large shapes but 83--87\% slow
  on B3; and B3 with K=256 behaves differently from B3 (slow on 16$\times$6, steady on 16$\times$8). ``NR of 6 or less
  is steady'' is a rule for B3 at K=384, not a general one.
  \item \textbf{No tile wins every shape.} For B3, 16$\times$6 reaches 6.71 / 6.77 GOPs (\texttt{jbpi2} / \texttt{jbpi})
  against 4.43 / 4.53 for the shipped 16$\times$15; but B2 is lower on 16$\times$6 (10.27 / 10.25 against
  12.59 / 12.56).
\end{enumerate}

\section{Hypothesis}
\textbf{Working hypothesis (not tested here).} The B packer reads NR columns of B at once, each a separate read stream.
The L1 prefetcher can follow up to six streams; at seven or more it backs off to about half its normal work (the pattern
\iss{41} and \iss{42} counted on 16$\times$15), the lines arrive late, and the rep slows. This is the suspicion of \iss{44}.

\emph{What it explains:} the threshold between six and seven at every MR (finding 1); why MR and kc do not matter
(finding 2); why the scalar packer, which reads the same addresses in the same order but one element at a time, goes
slow later and less (finding 3).

\emph{What it does not explain yet:} why compute slows as well as packing (finding 4); why 32$\times$7 is steady on shapes with
many rows and why B3 with K=256 is slow at 16$\times$6 (finding 5). So the stream count is at most one of the conditions.

\emph{How to test it:} \iss{44} (pack B one column at a time) predicts that a 15-column panel then behaves like a
6-column one. Counting prefetcher events on 16$\times$6 against 16$\times$7 (\texttt{-{}-raw-events x60-l2}) shows whether
the prefetch hits halve between them. \textbf{Nothing in this sweep counts prefetcher events;} the stream count is read from the
pattern in the data.

\section{Conclusion and impact}
\textbf{Conclusion.} Yes: the microkernel tile decides whether B3's slow state appears, and it does so through NR, the
number of B columns the packer reads at once. Packing also depends on the tile: in fast reps the vector packer needs
0.97--0.99 ms at 16$\times$6 against 0.66--0.76 ms at NR of 8 to 15, and 1.32 ms at 32$\times$6. The tile does \emph{not} explain
the reference design (\iss{46}): our 16$\times$14 is as slow as our 16$\times$15 (81--86\% against 89--92\% slow reps),
while the reference's 16$\times$14 is steady, so the difference between the two designs lies elsewhere (packers, buffer
placement or kernels).

\textbf{Impact.}
\begin{itemize}
  \item For small-M shapes such as B3, the shipped 16$\times$15 tile is the wrong choice: 16$\times$6 gives
  6.71 / 6.77 GOPs against 4.43 / 4.53 (and Reuse24 about 9.1 against 8.5--8.8).
  \item The same tile costs B2 (10.27 against 12.59 GOPs), and B3 with K=256 is slow on 16$\times$6. No single tile is best,
  so the tile has to be chosen per shape; that is the user's decision and out of scope here.
  \item The vector packer is where the tile matters most: with seven or more B streams it runs at about half the speed of the scalar
  packer (16$\times$8: 3.6--3.7 against 7.0--7.3 GOPs). If \iss{44} confirms the hypothesis, a change to the packer, not to the tile,
  could keep wide tiles steady.
\end{itemize}


\section{What the vector packers do for each tile}
The shipped packers are \texttt{pack\_RB\_v\_if} (A) and \texttt{pack\_CB\_v\_if} (B). The A packer loads MR/8 contiguous
vectors per k-step. The B packer loads ceil(NR/8) \emph{strided} vectors per k-step, each element from a different
column of B, so one pack walks \textbf{NR read streams in parallel}; NR is also the width of a row of the packed panel.
\iss{44} suspects that the prefetcher cannot follow that many streams. The cache model also changes kc with the tile.

\begin{center}\footnotesize\setlength{\tabcolsep}{3pt}
\begin{tabular}{@{}lcccccccc@{}}
\toprule
tile & kc & A vec.\ per k & B vec.\ per k & B streams & Bc row (B) & Bc stores over a line (\%) & live vregs & fringe / all tiles \\
\midrule
16$\times$6 & 256 & 2 & 1 & 6 & 24 & 25 & 14 & 0/256 \\
16$\times$7 & 256 & 2 & 1 & 7 & 28 & 38 & 16 & 4/220 \\
16$\times$8 & 256 & 2 & 1 & 8 & 32 & 0 & 18 & 0/192 \\
16$\times$14 & 128 & 2 & 2 & 14 & 56 & 31 & 30 & 4/112 \\
16$\times$15 & 128 & 2 & 2 & 15 & 60 & 41 & 32 & 4/104 \\
16$\times$16 & 128 & 2 & 2 & 16 & 64 & 0 & 34 & 0/96 \\
32$\times$6 & 128 & 4 & 1 & 6 & 24 & 25 & 28 & 0/128 \\
32$\times$7 & 128 & 4 & 1 & 7 & 28 & 38 & 32 & 2/110 \\
8$\times$8 & 256 & 1 & 1 & 8 & 32 & 0 & 9 & 0/384 \\
4$\times$4 & 384 & 1 & 1 & 4 & 16 & 0 & 5 & 0/1536 \\
\bottomrule
\end{tabular}
\end{center}
A tile of 16$\times$16 needs 34 live vector registers, over the limit of 32; it is the one known-bad reference in the
sweep. A tile at the edge of the block (a fringe tile) is computed through a temporary buffer, which belongs to compute
time and moves with the tile.

\section{B3 on every tile, vector packer}
GOPs by mean time, GOPs by median rep, and the share of slow reps (\%), averaged over the launches (4 on \texttt{jbpi2},
3 on \texttt{jbpi}). Green: no slow rep in any launch.

\begin{center}\footnotesize\setlength{\tabcolsep}{4pt}
\begin{tabular}{@{}lcccccccc@{}}
\toprule
 & & \multicolumn{3}{c}{\texttt{jbpi2}} & \multicolumn{3}{c}{\texttt{jbpi}} \\
\cmidrule(lr){3-5}\cmidrule(lr){6-8}
tile & kc & mean time & median rep & slow \% & mean time & median rep & slow \% \\
\midrule
4$\times$4 & 384 & \good{3.55} & 3.55 & \good{0} & \good{3.55} & 3.56 & \good{0} \\
4$\times$8 & 384 & 2.75 & 2.65 & \up{78} & 2.62 & 2.47 & \up{82} \\
4$\times$16 & 256 & 2.62 & 2.46 & \up{81} & 2.73 & 2.47 & \up{76} \\
8$\times$4 & 384 & \good{4.91} & 4.91 & \good{0} & \good{4.91} & 4.91 & \good{0} \\
8$\times$8 & 256 & 3.46 & 3.17 & \up{88} & 3.84 & 3.50 & \up{73} \\
8$\times$14 & 256 & 3.61 & 3.29 & \up{89} & 3.48 & 3.26 & \up{93} \\
8$\times$15 & 256 & 3.37 & 3.27 & \up{92} & 3.54 & 3.29 & \up{90} \\
8$\times$16 & 256 & 3.57 & 3.27 & \up{88} & 3.75 & 3.29 & \up{81} \\
8$\times$24 & 170 & 4.92 & 5.24 & \up{54} & 5.20 & 5.60 & \up{53} \\
8$\times$31 & 132 & 4.32 & 4.17 & \up{75} & 4.11 & 3.80 & \up{84} \\
12$\times$8 & 170 & 2.44 & 2.35 & \up{72} & 2.36 & 2.18 & \up{82} \\
16$\times$4 & 256 & \good{5.32} & 5.32 & \good{0} & \good{5.30} & 5.30 & \good{0} \\
16$\times$6 & 256 & \good{6.71} & 6.72 & \good{0} & \good{6.77} & 6.79 & \good{0} \\
16$\times$7 & 256 & 4.76 & 4.93 & \up{68} & 4.44 & 4.36 & \up{79} \\
16$\times$8 & 256 & 3.56 & 3.30 & \up{94} & 3.72 & 3.50 & \up{89} \\
16$\times$12 & 128 & 5.10 & 5.09 & \up{78} & 4.71 & 4.35 & \up{88} \\
16$\times$14 & 128 & 4.72 & 4.59 & \up{86} & 4.51 & 4.20 & \up{81} \\
16$\times$15 & 128 & 4.43 & 4.31 & \up{92} & 4.53 & 4.31 & \up{89} \\
16$\times$16 & 128 & 3.53 & 3.14 & \up{67} & 3.29 & 3.11 & 25 \\
24$\times$6 & 170 & \good{5.35} & 5.36 & \good{0} & \good{5.38} & 5.38 & \good{0} \\
24$\times$8 & 170 & 3.30 & 3.05 & \up{89} & 3.33 & 3.06 & \up{87} \\
24$\times$9 & 170 & 3.69 & 3.54 & \up{83} & 3.55 & 3.26 & \up{89} \\
32$\times$4 & 128 & \good{5.78} & 5.78 & \good{0} & \good{5.80} & 5.80 & \good{0} \\
32$\times$6 & 128 & \good{6.33} & 6.34 & \good{0} & \good{6.32} & 6.32 & \good{0} \\
32$\times$7 & 128 & 5.39 & 5.36 & \up{83} & 5.30 & 5.12 & \up{87} \\
\bottomrule
\end{tabular}
\end{center}

\paragraph{How to read it.} The seven green rows are exactly the tiles with NR of 4 or 6. Every row with NR of 7 or
more is slow in at least half of the reps, except 16$\times$16 on \texttt{jbpi} (25\%); the least slow of the others
is 8$\times$24 (53--54\%). The fastest tiles are the NR\,=\,6 ones: 16$\times$6 and 32$\times$6.

\section{B3 on the scalar packer}
Same measures, on the seven tiles of the scalar arm. The vector figures for the same tiles are in the table above.

\begin{center}\footnotesize\setlength{\tabcolsep}{4pt}
\begin{tabular}{@{}lcccccc@{}}
\toprule
 & \multicolumn{3}{c}{\texttt{jbpi2}} & \multicolumn{3}{c}{\texttt{jbpi}} \\
\cmidrule(lr){2-4}\cmidrule(lr){5-7}
tile & mean time & median rep & slow \% & mean time & median rep & slow \% \\
\midrule
4$\times$4 & \good{3.39} & 3.39 & \good{0} & \good{3.38} & 3.39 & \good{0} \\
8$\times$15 & 4.73 & 4.85 & 49 & 4.83 & 5.14 & 46 \\
16$\times$8 & 7.28 & 7.44 & 2 & 7.03 & 7.40 & 8 \\
16$\times$14 & 6.98 & 7.78 & 33 & 7.19 & 7.96 & 32 \\
16$\times$15 & 6.77 & 7.51 & 39 & 6.13 & 5.99 & \up{52} \\
16$\times$16 & 3.83 & 3.74 & \up{60} & 3.79 & 3.96 & \up{55} \\
32$\times$6 & \good{5.89} & 5.89 & \good{0} & \good{5.92} & 5.92 & \good{0} \\
\bottomrule
\end{tabular}
\end{center}
The scalar packer reads the same addresses in the same order, one element at a time. It is not free of the slow
state (16$\times$14 to 16$\times$16 and 8$\times$15 are 32--60\% slow), but it enters it less often, and the gap to
the vector packer is largest at NR\,=\,8, where the vector packer runs at about half the speed.

\section{The controls on the same tiles}
Vector packer, GOPs by mean time (\texttt{jbpi2} then \texttt{jbpi} in each pair); in brackets, the share of slow reps
when it is not zero. Green: no slow rep.

\begin{center}\footnotesize\setlength{\tabcolsep}{4pt}
\begin{tabular}{@{}lcccccccc@{}}
\toprule
 & \multicolumn{2}{c}{B3} & \multicolumn{2}{c}{B3\_K256} & \multicolumn{2}{c}{B2} & \multicolumn{2}{c}{Reuse24} \\
\cmidrule(lr){2-3}\cmidrule(lr){4-5}\cmidrule(lr){6-7}\cmidrule(lr){8-9}
tile & jbpi2 & jbpi & jbpi2 & jbpi & jbpi2 & jbpi & jbpi2 & jbpi \\
\midrule
16$\times$6 & \good{6.71} & \good{6.77} & 4.49~(58\%) & 4.96~(48\%) & \good{10.27} & \good{10.25} & \good{9.14} & \good{9.11} \\
16$\times$7 & 4.76~(68\%) & 4.44~(79\%) & 4.09~(81\%) & 4.09~(79\%) & \good{10.76} & \good{10.73} & \good{9.53} & \good{9.54} \\
16$\times$8 & 3.56~(94\%) & 3.72~(89\%) & \good{6.74} & \good{6.76} & \good{10.24} & \good{10.33} & \good{8.94} & \good{8.64} \\
16$\times$14 & 4.72~(86\%) & 4.51~(81\%) & 4.44~(93\%) & 4.52~(93\%) & \good{11.87} & \good{11.84} & \good{8.29} & \good{8.31} \\
16$\times$15 & 4.43~(92\%) & 4.53~(89\%) & 4.93~(88\%) & 4.76~(94\%) & \good{12.59} & \good{12.56} & 8.45~(17\%) & 8.79~(4\%) \\
16$\times$16 & 3.53~(67\%) & 3.29~(25\%) & 3.45~(88\%) & 3.73~(77\%) & \good{7.95} & \good{7.95} & \good{4.93} & \good{4.89} \\
32$\times$6 & \good{6.33} & \good{6.32} & 5.54~(30\%) & 5.54~(19\%) & 10.65~(0\%) & \good{10.62} & \good{9.37} & \good{9.35} \\
32$\times$7 & 5.39~(83\%) & 5.30~(87\%) & 5.93~(62\%) & 6.01~(59\%) & \good{11.80} & \good{11.70} & \good{9.81} & \good{9.81} \\
24$\times$6 & \good{5.35} & \good{5.38} & 3.56~(85\%) & 3.57~(81\%) & \good{9.33} & \good{9.30} & 6.36~(14\%) & 6.29~(69\%) \\
8$\times$8 & 3.46~(88\%) & 3.84~(73\%) & 3.54~(89\%) & 3.39~(95\%) & \good{8.14} & \good{8.14} & \good{7.55} & \good{7.52} \\
4$\times$4 & \good{3.55} & \good{3.55} & \good{3.93} & \good{3.91} & \good{4.19} & \good{4.20} & \good{4.21} & \good{4.22} \\
\bottomrule
\end{tabular}
\end{center}
B2 and Reuse24 are steady on almost every tile, so on the very launches where B3 is 90\% slow, the board is not the
cause. B3 with K=256 follows a different rule: it is steady on 16$\times$8 and slow on 16$\times$6. No tile is best for all four
shapes: 16$\times$6 is best for B3 and Reuse24 and costs B2 about 2.3 GOPs; 16$\times$15 is best for B2 and costs B3
about 2.2 GOPs.

\section{Packing and compute apart}
Instrumented build, B3; these are debug numbers, not speed. Mean milliseconds per rep, pooled over the launches: fast
reps (pack / compute), slow reps (pack / compute), and pack's share of pack plus compute (\%).

\begin{center}\footnotesize\setlength{\tabcolsep}{3pt}
\begin{tabular}{@{}lcccccc@{}}
\toprule
 & \multicolumn{3}{c}{\texttt{jbpi2}} & \multicolumn{3}{c}{\texttt{jbpi}} \\
\cmidrule(lr){2-4}\cmidrule(lr){5-7}
tile, packer & fast reps & slow reps & pack \% & fast reps & slow reps & pack \% \\
\midrule
16$\times$6 vector & 0.97 / 1.85 & none & 34 & 0.99 / 1.86 & none & 35 \\
32$\times$6 vector & 1.32 / 1.73 & none & 43 & 1.32 / 1.74 & none & 43 \\
8$\times$8 vector & 0.79 / 2.27 & 1.70 / 3.81 & 30 & 0.83 / 2.29 & 1.69 / 3.76 & 31 \\
16$\times$8 vector & 0.76 / 1.86 & 1.69 / 3.32 & 33 & 0.76 / 1.88 & 1.69 / 3.29 & 33 \\
16$\times$14 vector & 0.72 / 1.89 & 1.52 / 2.74 & 33 & 0.66 / 1.80 & 1.43 / 2.66 & 32 \\
16$\times$15 vector & 0.71 / 1.80 & 1.48 / 2.68 & 34 & 0.73 / 1.79 & 1.46 / 2.67 & 34 \\
16$\times$8 scalar & 0.77 / 1.79 & (2 reps) & 30 & 0.80 / 1.79 & 1.50 / 2.68 & 31 \\
16$\times$15 scalar & 0.71 / 1.51 & 1.13 / 1.82 & 33 & 0.72 / 1.51 & 1.39 / 2.00 & 35 \\
32$\times$6 scalar & 1.50 / 1.73 & none & 46 & 1.50 / 1.74 & none & 46 \\
\bottomrule
\end{tabular}
\end{center}
\begin{itemize}
  \item In a slow rep \textbf{both halves lengthen}, on every tile that has slow reps. Tiles that never go slow have no
  slow reps in either half.
  \item In fast reps the pack time depends on the tile: 0.97--0.99 ms on 16$\times$6 against 0.66--0.76 ms on
  16$\times$8 to 16$\times$15, and 1.32 ms on 32$\times$6. The data moved is the same for B3 on every tile, so this is
  cost per vector operation and per panel; the sweep does not say which.
  \item The instrumented build is not the same experiment as the speed runs. 32$\times$7 has no slow reps on
  \texttt{jbpi2} and 578 of 2997 on \texttt{jbpi} there, against 83--87\% slow reps in the uninstrumented build on
  both boards. Only the fast-against-slow split is read from it.
\end{itemize}

\section{Conditions, limits and next steps}
\paragraph{Conditions.} All 26 run folders record \texttt{other users : none} at the start of the run, no series wrote
a load warning, and the processes read before and after every run list no other user's job, so there is no
``other users present'' note to attach to these numbers. The seven unpinned cores of each board were at 0--1\% busy
in every per-tile timing window but two, both in the instrumented build: 14\% on core 7 during the 8$\times$14 tile on
\texttt{jbpi}, and 2--4\% on cores 3 and 5 during the 8$\times$31 tile on \texttt{jbpi2}. Neither affects a speed
number, and we do not know whether it came from outside the benchmark. A \texttt{psiwins} job of 1.6--1.8 cores had been on both
boards at 12:19; it was gone by 14:04, and the short smoke test made while it ran was discarded. The two boards ran
simultaneously (separate hardware; only the NFS server and the jump host are shared). The runs were made in the
afternoon and evening, not at midnight; no other user's job appeared, so a quiet-window repeat was not needed.

\paragraph{Limits.}
\begin{itemize}
  \item Three or four launches per case. The \texttt{jbpi2} vector arm includes one launch from a series that was
  stopped and restarted when the work was split across the boards.
  \item The cause is inferred, not measured: nothing here counts prefetcher events. That the threshold lies between six
  and seven read streams is the pattern in the data; \iss{44} is the direct test.
  \item kc moves with the tile. The steady and the slow tiles share kc values, so kc does not explain the split, but
  the tiles are not a controlled change of NR alone at fixed kc.
  \item B3 only is analysed in detail. The other shapes are in the full report (Appendix~A).
\end{itemize}

\paragraph{Next steps.}
\begin{itemize}
  \item Choosing a tile for small-M shapes is a decision on this evidence; the shapes of the real models need the same
  sweep first.
  \item Run \iss{44} (pack B one column at a time): this brief predicts that a 15-column panel should then behave like a
  6-column one.
  \item Count prefetcher events on 16$\times$6 against 16$\times$7 (\texttt{-{}-raw-events x60-l2}, \iss{42}), the
  smallest pair that flips the state. \textbf{Now \iss{50}}: widened to the two B3\_K256 tiles that flip the other way
  (16$\times$6 and 16$\times$8), script \texttt{run\_issue50.sh}, not yet run (needs board access).
  \item Opened as \iss{49}: whether the B3\_K256 flip follows \texttt{ldb} rather than NR (the premise \iss{44} rests
  on). Its desk table (no board) already found that none of stream count, page/offset geometry, kc/K-blocks or a
  simple MR threshold alone splits every steady tile from every slow one across both K=384 and K=256 --- the decisive
  test (padding \texttt{ldb} directly, \texttt{run\_issue49.sh}) is the next step, also not yet run.
\end{itemize}

\section*{Appendix A: logs and code}
\raggedright\emergencystretch=3em
All paths are under \texttt{tests/pure\_xdsl\_blis\_rvv/}; all are committed (commit \texttt{476d843}).
\begin{itemize}
  \item Series: \path{logs/issue48_2026-09-30_124654} (stopped after 2 runs: \path{run_2026-09-30_134329},
  \path{run_2026-09-30_135936}); \path{logs/issue48_2026-09-30_140902} (\texttt{jbpi2}, 12 runs);
  \path{logs/issue48_2026-09-30_140902_jbpi} (\texttt{jbpi}, 12 runs). Each series' \texttt{runs.tsv} lists its run folders.
  \item Full report, every shape and tile: \path{logs/issue48_2026-09-30_report.md}
  (\texttt{python}~\path{issue48_report.py} on the three series folders).
  \item Code: \path{pack_dims_table.py} (the per-tile facts, tested against the builders),
  \path{config_issue48_sweep.json} and \path{config_issue48_sweep_scalar.json} (the tiles),
  \path{run_issue48.sh} (the series; \texttt{-{}-workspace} for one series per board), \path{issue48_report.py}.
  \item Write-up and the written answer: \path{.scratch/issues/0032-b3-microkernel-dimension-sweep-pack-and-compute.md}, \S7--\S8.
\end{itemize}

\end{document}
